\documentclass[11pt]{article}
\usepackage[utf8]{inputenc}
\usepackage{amsmath}
\usepackage{booktabs}
\usepackage[a4paper,margin=2.5cm]{geometry}

\newcommand{\iso}[2]{\textsuperscript{#1}#2}
\newcommand{\Rc}{\ensuremath{R_{50}}}
\newcommand{\Roc}{\ensuremath{R_{80}}}

\begin{document}

\begin{center}
  {\large\textbf{Calibration systematics from cross sections and washout:}}\\[2pt]
  {\large\textbf{plan}}
\end{center}
\medskip

\noindent
Plan for quantifying $\sigma_{\Delta R}$, the calibration term of
$\sigma_{\Roc} = \sigma_{\Delta R} \oplus \sigma_{\Rc}$ (\texttt{cbs.tex}
Eq.~8), from two sources: nuclear production cross sections and the Mizuno
washout parameters. \texttt{cbs.tex} measures $\sigma_{\Rc}$ (0.101~mm BGO ring,
0.146~mm CsI ring, washed, del120) and declines $\sigma_{\Delta R}$.

\subsection*{1. Blind rule}

Monte Carlo truth (\texttt{isotope}, \texttt{z0\_mm}, production energy) is used
\emph{only} to build the perturbed pseudo-data. The analysis is the existing
pipeline unchanged: whole-plane depth profile from the reconstructed image,
bounded erfc fit, \Rc. No isotope label, no true production point, no
perturbation parameter enters the analysis. Scan observables --- reconstructed
position, energy, and decay time \texttt{t\_decay\_s} --- are admissible, but
Stage~0 uses none beyond the profile.

The calibration $\Delta R = \Roc - \Rc$ is computed once from the nominal Monte
Carlo and held fixed. Nuclear production cross sections do not move \Roc\ (the
dose edge is set by electronic stopping power), so the bias is

\begin{equation}
  b = \Rc(\text{perturbed}) - \Rc(\text{nominal}),
\end{equation}

read with the frozen calibration. This is the quantity to tabulate.

\subsection*{2. What moves \Rc}

Write the profile as a mixture of per-isotope depth shapes,
\begin{equation}
  P(z) = \sum_i w_i\,\hat P_i(z),
  \qquad \textstyle\int \hat P_i = 1 .
\end{equation}
Uniform washout scales $w_i$ only. Cross sections scale $w_i$ (channel yields)
and distort $\hat P_i(z)$ (energy dependence of $\sigma_i(E)$ near threshold).

\paragraph{Mixture lever, calibrated.} Washout multiplies the mix by
$g_i$ = 0.448 (\iso{15}{O}), 0.376 (\iso{11}{C}), raising the \iso{15}{O} window
fraction from 0.82 to 0.844, and shifts the edge by the measured $+0.218$~mm.
Hence
\begin{equation}
  \frac{\mathrm{d}\Rc}{\mathrm{d}f_{^{15}\mathrm{O}}} \simeq 9\text{--}10~\text{mm per unit fraction},
\end{equation}
consistent with the delayed-start calibration walk ($-2.9$~mm as \iso{15}{O}
drains by $t_{\mathrm{start}} = 300$~s). \textbf{A 1\% error in the \iso{15}{O}
fraction costs $\simeq 0.09$~mm}, comparable to $\sigma_{\Rc}$.

\paragraph{Shape lever, not calibrated away.} A $z$-dependent distortion of one
isotope's own production profile changes the shape of the distal fall-off. No
scalar absorbs it. Structurally identical to the spatial-non-uniformity washout
item (\texttt{md/isotope-washout.md}), which is the one known non-calibratable
washout bias. Cross-section libraries disagree most near threshold and below
$\sim$30~MeV; the distal few mm of activity is produced by protons in that
energy range, and in a SOBP the distal fall-off is formed by the
highest-energy layer alone. Expected to dominate; magnitude unmeasured.

\subsection*{3. Perturbations}

Each is a per-event weight on the existing pooled shards.

\begin{table}[ht]
  \centering
  \begin{tabular}{llll}
    \toprule
    & Perturbation & Applied as & Machinery \\
    \midrule
    P1 & Mizuno parameters & $g_i(M_k,T_k)$, Bernoulli keep & exists \\
    P2 & Cross-section yields & $r_i$ constant per isotope & exists ($g_i$ path) \\
    P3 & Cross-section shape & $r_i(z_0)$ per event & new, small \\
    \bottomrule
  \end{tabular}
\end{table}

\noindent
\textbf{P1.} Draw $(M_k, T_k)$ from the Mizuno uncertainties
(\texttt{config/washout\_brain.toml}), recompute $g_i$, keep with probability
$p_{\mathrm{dose}}\,g_i$. \texttt{tools/washout.py::mc\_band} already does this
at truth level and returns $\pm 0.020$~mm --- for the old window and the old
unbounded fit. Recompute at the reference protocol with the bounded fit.

\noindent
\textbf{P2.} Scale each channel yield by $r_i$ and keep with probability
$\propto r_i$. Identical code path to $g_i$.

\noindent
\textbf{P3.} Weight by $r_i(z_0)$ using the stored annihilation depth. Two ways
to set $r_i(z)$:
\begin{enumerate}
  \item[(a)] \emph{Parameterised.} Apply a shift or stretch to $\hat P_i(z)$ and
    report sensitivity per mm of production-profile shift. Model-independent;
    decouples what the method does from how bad the input is.
  \item[(b)] \emph{Physical.} Fold a cross-section ratio $r_i(E)$ with the proton
    fluence. In the distal region the residual-range relation gives $E(z)$
    directly; over the rest of the fit window it is approximate.
\end{enumerate}
Start with (a); (b) only if a defensible $r_i(E)$ is in hand.

\subsection*{4. Inputs}

Available: per-LOR \texttt{isotope}, \texttt{z0\_mm}, \texttt{t\_decay\_s}
(\texttt{dev/reference/SCHEMA.md}); per-isotope truth depth columns
(\texttt{truth/activity\_profile\_fast.csv}); Mizuno uncertainties.

\textbf{Gating item:} the per-channel cross-section spread for
\iso{16}{O}$(p,x)$\iso{15}{O}, \iso{12}{C}$(p,pn)$\iso{11}{C},
\iso{16}{O}$(p,x)$\iso{11}{C}, \iso{16}{O}$(p,x)$\iso{13}{N}. Source:
Espa\~na et al.\ 2011 (already cited in \texttt{cbs.tex}), which compares
datasets and their effect on predicted PET range. We cannot run FLUKA --- that
is a second transport chain upstream of this repository. Using the published
inter-dataset spread is the defensible substitute; claiming a FLUKA counterpart
is not.

\subsection*{5. Procedure --- Stage 0 (blind, no correction)}

\begin{enumerate}
  \item Fix the calibration from the nominal Monte Carlo at the reference
    protocol (del120, 1~Gy).
  \item For each perturbation P1--P3, and for a grid of its parameter, reweight
    the pool, refit \Rc\ blind, record $b$.
  \item Report $b$ per source and the combined band.
\end{enumerate}

Run at \textbf{origins level first} --- reweight the truth depth columns, refit,
no reconstruction. Bias transfers scanner-independently (washout: truth
$+0.218$~mm vs detected $+0.20$--$0.25$~mm), so this is nearly free and answers
the physics question. Then confirm the total on one reconstructed arm (BGO ring)
through \texttt{drivers/statistical\_procedure\_jobs.jl}.

\subsection*{6. Later stages, conditional on Stage 0}

\noindent
\textbf{Stage 1.} If $b$ is large, ask how much a calibration constant absorbs
--- i.e.\ how much of $b$ is a rigid shift versus a shape change.

\noindent
\textbf{Stage 2.} In-situ constraint from the time structure. $P(z,t)$
decomposes by half-life; over the reference window \iso{15}{O} decays by 0.18
while \iso{11}{C} decays by 0.84, a 4.6$\times$ contrast. Unfolding gives
$w_i \hat P_i(z)$ from scan observables only, which is legal under the blind
rule (\texttt{t\_decay\_s} is acquired). This converts both systematics from
priors into measured quantities, at a statistical cost bounded below by the
per-isotope $\sigma_{\Rc}$ already measured
($k = \sigma\sqrt{N} \approx 236$ for \iso{15}{O}, 319 for \iso{11}{C}, ring).
Cost: one MLEM reconstruction per time bin per realisation; the
\texttt{--tend} window sub-cut already performs the required cut. Subsumes the
deferred composite-erfc item (\texttt{md/pending.md}).

\subsection*{7. Deliverables}

A table of $b$ [mm] per perturbation at the reference protocol, origins level
and one reconstructed arm; the combined $\sigma_{\Delta R}$ band from these two
sources; and the sensitivity coefficients
$\mathrm{d}\Rc/\mathrm{d}f_{^{15}\mathrm{O}}$ and $\mathrm{d}\Rc/\mathrm{d}(\text{profile shift})$.

If it enters \texttt{cbs.tex}: two sources out of many, on a uniform phantom
with no CT, no tissue heterogeneity, no anatomy, no positioning error.
Heterogeneity is likely the largest real contributor and is absent by
construction. Frame as demonstrating the method for isolating these
systematics, not as a systematic budget.

\end{document}
